\begin{enumerate}[label=\thesubsection.\arabic*, ref=\thesubsection.\theenumi]
	
\item Find $\mydet{\vec{AB}}$, if
\begin{align}
	\vec{A} = \myvec{0 & -1 \\ 0 & 2}, \quad \vec{B} = \myvec{3 & 5 \\ 0 & 0}
\end{align}
\hfill (12~2019)

\solution
From \eqref{eq:CBSE12-2019-det_mult}:
\begin{align}
	\mydet{\vec{AB}} = \mydet{\vec{A}} \mydet{\vec{B}}
\end{align}
From \eqref{eq:CBSE12-2019-det_2x2}, evaluating the individual determinants:
\begin{align}
	\mydet{\vec{A}} &= \mydet{\myvec{0 & -1 \\ 0 & 2}} = \brak{0}\brak{2} - \brak{-1}\brak{0} = 0 \\
	\mydet{\vec{B}} &= \mydet{\myvec{3 & 5 \\ 0 & 0}} = \brak{3}\brak{0} - \brak{5}\brak{0} = 0
\end{align}
Therefore:
\begin{align}
	\mydet{\vec{AB}} = 0 \times 0 = 0
\end{align}


\item Differentiate $e^{3x}$ with respect to $x$.
\hfill (12~2019)

\solution
Let the function be $y = e^{3x}$.
From \eqref{eq:CBSE12-2019-diff_exp_linear}, differentiating with respect to $x$ using the chain rule \eqref{eq:CBSE12-2019-diff_exp_chain}:
\begin{align}
	\frac{dy}{dx} = \frac{d}{dx}\sbrak{e^{3x}} = 3e^{3x}
\end{align}


\item Find the order and degree (if defined) of the differential equation
\begin{align}
	\frac{d^2y}{dx^2} + x\brak{\frac{dy}{dx}}^2 = 2x^2\ln\brak{\frac{d^2y}{dx^2}} \label{eq:EE26BTECH11223-12-q3-de}
\end{align}
\hfill (12~2019)

\solution
Following the criteria in \eqref{eq:CBSE12-2019-dc_def}:
\begin{itemize}
	\item The highest-order derivative appearing in \eqref{eq:EE26BTECH11223-12-q3-de} is $\dfrac{d^2y}{dx^2}$. Hence, $\text{Order} = 2$.
	\item The highest-order derivative $\dfrac{d^2y}{dx^2}$ appears inside a transcendental logarithmic function $\ln\brak{\cdot}$, meaning the differential equation cannot be expressed as a polynomial in its derivatives. Therefore, the \textbf{Degree is not defined}.
\end{itemize}

\item Find the direction cosines of a line which makes equal angles with the coordinate axes.
\hfill (12~2019)

\solution
Let the line make equal angles $\alpha$ with the coordinate axes. From \eqref{eq:CBSE12-2019-dc_def} and \eqref{eq:CBSE12-2019-dc_sum_sq}:
\begin{align}
	l^2 + m^2 + n^2 = 1 \implies 3l^2 = 1 \implies l = \pm\frac{1}{\sqrt{3}}
\end{align}
From \eqref{eq:CBSE12-2019-dc_equal}, since $l = m = n$, the direction cosines are:
\begin{align}
	\brak{\pm\frac{1}{\sqrt{3}}, \, \pm\frac{1}{\sqrt{3}}, \, \pm\frac{1}{\sqrt{3}}}
\end{align}

\item A line passes through the point with position vector $\vec{a} = \myvec{2 \\ -1 \\ 4}$ and is in the direction of the vector $\vec{b} = \myvec{1 \\ 1 \\ -2}$. Find the equation of the line in Cartesian form.
\hfill (12~2019)

\solution
From \eqref{eq:CBSE12-2019-line_cart}, with point $\vec{a} = \myvec{x_1 \\ y_1 \\ z_1} = \myvec{2 \\ -1 \\ 4}$ and direction vector $\vec{b} = \myvec{a \\ b \\ c} = \myvec{1 \\ 1 \\ -2}$:
\begin{align}
	\frac{x - 2}{1} = \frac{y - \brak{-1}}{1} = \frac{z - 4}{-2}
\end{align}
Simplifying:
\begin{align}
	\frac{x - 2}{1} = \frac{y + 1}{1} = \frac{z - 4}{-2}
\end{align}
The following code gives the plot
\begin{lstlisting}[frame=single, basicstyle=\ttfamily\small]
	codes/12/q5.py
\end{lstlisting}

\item Find $\displaystyle\int \brak{3 - 2x - x^2}\,dx$.
\hfill (12~2019)

\solution
Let $I = \int \brak{3 - 2x - x^2}\,dx$.
Applying linearity and the power integration rule \eqref{eq:CBSE12-2019-int_power}:
\begin{align}
	I &= \int 3\,dx - \int 2x\,dx - \int x^2\,dx \\
	&= 3x - 2\brak{\frac{x^2}{2}} - \frac{x^3}{3} + C \\
	&= 3x - x^2 - \frac{x^3}{3} + C
\end{align}


\item If $\vec{A} = \myvec{p & 2 \\ 2 & p}$ and $\mydet{\vec{A}^3} = 125$, find the values of $p$.
\hfill (12~2019)

\solution
From \eqref{eq:CBSE12-2019-det_2x2}, evaluating the determinant of $\vec{A}$:
\begin{align}
	\mydet{\vec{A}} = \mydet{\myvec{p & 2 \\ 2 & p}} = p^2 - 4
\end{align}
From \eqref{eq:CBSE12-2019-det_power}:
\begin{align}
	\mydet{\vec{A}^3} = \brak{\mydet{\vec{A}}}^3 = \brak{p^2 - 4}^3
\end{align}
Given that $\mydet{\vec{A}^3} = 125$:
\begin{align}
	\brak{p^2 - 4}^3 = 5^3 \implies p^2 - 4 = 5 \implies p^2 = 9 \implies p = \pm 3
\end{align}
Thus, the values of $p$ are $3$ and $-3$.


\item Examine whether the operation $*$ defined on the set of all real numbers by $a*b = a^2 + b^2$ is a binary operation or not, and if it is a binary operation, find whether it is associative or not.
\hfill (12~2019)

\solution
From \eqref{eq:CBSE12-2019-binary_closure}, for any $a, b \in \mathbb{R}$, $a^2 \in \mathbb{R}$ and $b^2 \in \mathbb{R}$. Since the sum of real numbers is real, $a^2 + b^2 \in \mathbb{R}$. Hence, $*$ is a well-defined binary operation on $\mathbb{R}$.
To test associativity from \eqref{eq:CBSE12-2019-associative}:
\begin{align}
	\text{LHS:} \quad \brak{a * b} * c &= \brak{a^2 + b^2} * c = \brak{a^2 + b^2}^2 + c^2 \\
	\text{RHS:} \quad a * \brak{b * c} &= a * \brak{b^2 + c^2} = a^2 + \brak{b^2 + c^2}^2
\end{align}
Choosing counterexample $a = 1, b = 2, c = 3 \in \mathbb{R}$:
\begin{align}
	\brak{1 * 2} * 3 &= \brak{1^2 + 2^2} * 3 = 5 * 3 = 5^2 + 3^2 = 25 + 9 = 34 \\
	1 * \brak{2 * 3} &= 1 * \brak{2^2 + 3^2} = 1 * 13 = 1^2 + 13^2 = 1 + 169 = 170
\end{align}
Since $34 \neq 170$, the operation $*$ is **not associative**.

\item If $\norm{\vec{a}} = 2$, $\norm{\vec{b}} = 7$, and $\vec{a} \times \vec{b} = \myvec{3 \\ 2 \\ 6}$, find the angle between $\vec{a}$ and $\vec{b}$.
\hfill (12~2019)

\solution
Let $\theta$ be the angle between $\vec{a}$ and $\vec{b}$.
Evaluating the magnitude of $\vec{a} \times \vec{b}$:
\begin{align}
	\norm{\vec{a} \times \vec{b}} = \sqrt{3^2 + 2^2 + 6^2} = \sqrt{9 + 4 + 36} = \sqrt{49} = 7
\end{align}
From \eqref{eq:CBSE12-2019-cross_norm} and \eqref{eq:CBSE12-2019-sin_theta}:
\begin{align}
	\sin\theta = \frac{\norm{\vec{a} \times \vec{b}}}{\norm{\vec{a}}\norm{\vec{b}}} = \frac{7}{2 \times 7} = \frac{1}{2}
\end{align}
Since $\theta \in [0, \pi]$:
\begin{align}
	\theta = \frac{\pi}{6} \quad \text{or} \quad \theta = \frac{5\pi}{6}
\end{align}


\item Find the volume of a cuboid whose edges are given by $\vec{a} = \myvec{-3 \\ 7 \\ 5}$, $\vec{b} = \myvec{-5 \\ 7 \\ -3}$, and $\vec{c} = \myvec{7 \\ -5 \\ -3}$.
\hfill (12~2019)

\solution
From \eqref{eq:CBSE12-2019-scalar_triple_vol}, the volume $V$ of a parallelepiped (cuboid) is the absolute value of the scalar triple product:
\begin{align}
	V = \abs{\mydet{\vec{a} & \vec{b} & \vec{c}}} = \abs{\mydet{-3 & -5 & 7 \\ 7 & 7 & -5 \\ 5 & -3 & -3}}
\end{align}
Expanding the determinant:
\begin{align}
	\mydet{\vec{a} & \vec{b} & \vec{c}} &= -3\mydet{7 & -5 \\ -3 & -3} - \brak{-5}\mydet{7 & -5 \\ 5 & -3} + 7\mydet{7 & 7 \\ 5 & -3} \\
	&= -3\brak{-21 - 15} + 5\brak{-21 - \brak{-25}} + 7\brak{-21 - 35} \\
	&= -3\brak{-36} + 5\brak{4} + 7\brak{-56} = 108 + 20 - 392 = -264
\end{align}
Taking the absolute value:
\begin{align}
	V = \abs{-264} = 264\text{ cubic units}
\end{align}
The following code verifies the volume:
\begin{lstlisting}[frame=single, basicstyle=\ttfamily\small]
	codes/12/q10.py
\end{lstlisting}


\item A coin is tossed $5$ times. What is the probability of getting (i) $3$ heads, (ii) at most $3$ heads?
\hfill (12~2019)

\solution
Let $X$ denote the number of heads obtained in $n = 5$ independent coin tosses. This represents a sequence of Bernoulli trials with success probability $p = \frac{1}{2}$ and failure probability $q = 1 - p = \frac{1}{2}$.
From the binomial distribution formula \eqref{eq:CBSE12-2019-binomial_pmf}:
\begin{align}
	\pr{X = k} = \binom{n}{k} p^k q^{n-k} = \binom{5}{k} \brak{\frac{1}{2}}^k \brak{\frac{1}{2}}^{5-k} = \frac{\binom{5}{k}}{32}
\end{align}

(i) Probability of getting exactly $3$ heads ($k = 3$):
\begin{align}
	\pr{X = 3} = \frac{\binom{5}{3}}{32} = \frac{10}{32} = \frac{5}{16}
\end{align}

(ii) Probability of getting at most $3$ heads ($\pr{X \le 3}$):
Using the complement relation from \eqref{eq:CBSE12-2019-binomial_cdf}:
\begin{align}
	\pr{X \le 3} &= 1 - \sbrak{\pr{X = 4} + \pr{X = 5}} \\
	&= 1 - \left(\frac{\binom{5}{4}}{32} + \frac{\binom{5}{5}}{32}\right) = 1 - \frac{5 + 1}{32} = 1 - \frac{6}{32} = \frac{26}{32} = \frac{13}{16}
\end{align}


\item Find the probability distribution of $X$, the number of heads in a simultaneous toss of two coins.
\hfill (12~2019)

\solution
When two fair coins are tossed simultaneously, the sample space $S$ is:
\begin{align}
	S = \{\text{HH}, \, \text{HT}, \, \text{TH}, \, \text{TT}\}
\end{align}
Let $X$ be the random variable representing the number of heads. The possible values of $X$ are $0, 1, 2$.
From \eqref{eq:CBSE12-2019-prob_dist_axioms}:
\begin{align}
	\pr{X = 0} &= \frac{n(\{\text{TT}\})}{n(S)} = \frac{1}{4} \\
	\pr{X = 1} &= \frac{n(\{\text{HT}, \text{TH}\})}{n(S)} = \frac{2}{4} = \frac{1}{2} \\
	\pr{X = 2} &= \frac{n(\{\text{HH}\})}{n(S)} = \frac{1}{4}
\end{align}
Verifying the distribution sum:
\begin{align}
	\sum_{x=0}^{2} \pr{X = x} = \frac{1}{4} + \frac{1}{2} + \frac{1}{4} = 1
\end{align}


\item Find $\displaystyle\int \frac{\sin^3 x + \cos^3 x}{\sin^2 x \cos^2 x}\,dx$.
\hfill (12~2019)

\solution
Let $I = \int \frac{\sin^3 x + \cos^3 x}{\sin^2 x \cos^2 x}\,dx$.
Splitting the fraction into two terms:
\begin{align}
	\frac{\sin^3 x + \cos^3 x}{\sin^2 x \cos^2 x} &= \frac{\sin^3 x}{\sin^2 x \cos^2 x} + \frac{\cos^3 x}{\sin^2 x \cos^2 x} \\
	&= \frac{\sin x}{\cos^2 x} + \frac{\cos x}{\sin^2 x} \\
	&= \brak{\frac{\sin x}{\cos x}}\brak{\frac{1}{\cos x}} + \brak{\frac{\cos x}{\sin x}}\brak{\frac{1}{\sin x}} \\
	&= \sec x \tan x + \csc x \cot x
\end{align}
Integrating term-by-term using \eqref{eq:CBSE12-2019-int_sec_tan} and \eqref{eq:CBSE12-2019-int_csc_cot}:
\begin{align}
	I &= \int \brak{\sec x \tan x + \csc x \cot x}\,dx \\
	&= \int \sec x \tan x\,dx + \int \csc x \cot x\,dx \\
	&= \sec x - \csc x + C
\end{align}


\item Find $\displaystyle\int \frac{x - 3}{\brak{x - 1}^3} e^x\,dx$.
\hfill (12~2019)

\solution
Let $I = \int \frac{x - 3}{\brak{x - 1}^3} e^x\,dx$.
Expressing the numerator in terms of $\brak{x - 1}$:
\begin{align}
	\frac{x - 3}{\brak{x - 1}^3} = \frac{\brak{x - 1} - 2}{\brak{x - 1}^3} = \frac{1}{\brak{x - 1}^2} - \frac{2}{\brak{x - 1}^3}
\end{align}
Let $f(x) = \frac{1}{\brak{x - 1}^2} = \brak{x - 1}^{-2}$. Its derivative is:
\begin{align}
	f'(x) = -2\brak{x - 1}^{-3} = -\frac{2}{\brak{x - 1}^3}
\end{align}
Using the standard exponential integral formula \eqref{eq:CBSE12-2019-exp_f_plus_fprime}:
\begin{align}
	\int e^x \sbrak{f(x) + f'(x)}\,dx = e^x f(x) + C
\end{align}
Substituting $f(x)$:
\begin{align}
	I = \int e^x \sbrak{\frac{1}{\brak{x - 1}^2} - \frac{2}{\brak{x - 1}^3}}\,dx = \frac{e^x}{\brak{x - 1}^2} + C
\end{align}


\item If $\pr{\text{not } A} = 0.7$, $\pr{B} = 0.7$, and $\pr{B \mid A} = 0.5$, find $\pr{A \mid B}$.
\hfill (12~2019)

\solution
From the complement rule, $\pr{A} = 1 - \pr{\text{not } A} = 1 - 0.7 = 0.3$.
From the definition of conditional probability \eqref{eq:CBSE12-2019-cond_prob}:
\begin{align}
	\pr{B \mid A} = \frac{\pr{A \cap B}}{\pr{A}} \implies 0.5 = \frac{\pr{A \cap B}}{0.3} \implies \pr{A \cap B} = 0.15
\end{align}
Now, computing $\pr{A \mid B}$ via \eqref{eq:CBSE12-2019-cond_prob}:
\begin{align}
	\pr{A \mid B} = \frac{\pr{A \cap B}}{\pr{B}} = \frac{0.15}{0.7} = \frac{15}{70} = \frac{3}{14} \approx 0.2143
\end{align}


\item Find the general solution of the differential equation $\dfrac{dy}{dx} = e^{x+y}$.
\hfill (12~2019)

\solution
Using index laws, the differential equation can be separated:
\begin{align}
	\frac{dy}{dx} = e^x e^y \implies e^{-y}\,dy = e^x\,dx
\end{align}
Integrating both sides using \eqref{eq:CBSE12-2019-sep_vars}:
\begin{align}
	\int e^{-y}\,dy = \int e^x\,dx \implies -e^{-y} = e^x + C_1
\end{align}
Rearranging terms:
\begin{align}
	e^x + e^{-y} = C \quad \text{where } C = -C_1
\end{align}


\item Find the value of $\sin\left(\cos^{-1}\frac{4}{5} + \tan^{-1}\frac{2}{3}\right)$.
\hfill (12~2019)

\solution
From \eqref{eq:CBSE12-2019-arccos_to_arctan}, converting $\cos^{-1}\brak{\frac{4}{5}}$ to $\tan^{-1}$:
\begin{align}
	\cos^{-1}\frac{4}{5} = \tan^{-1}\left(\frac{\sqrt{1 - \brak{4/5}^2}}{4/5}\right) = \tan^{-1}\left(\frac{3/5}{4/5}\right) = \tan^{-1}\frac{3}{4}
\end{align}
Using the addition formula \eqref{eq:CBSE12-2019-arctan_add} for $\tan^{-1}\frac{3}{4} + \tan^{-1}\frac{2}{3}$ (since $uv = \frac{3}{4}\cdot\frac{2}{3} = \frac{1}{2} < 1$):
\begin{align}
	\tan^{-1}\frac{3}{4} + \tan^{-1}\frac{2}{3} &= \tan^{-1}\left(\frac{\frac{3}{4} + \frac{2}{3}}{1 - \frac{3}{4} \cdot \frac{2}{3}}\right) \\
	&= \tan^{-1}\left(\frac{\frac{17}{12}}{1 - \frac{6}{12}}\right) = \tan^{-1}\left(\frac{17/12}{6/12}\right) = \tan^{-1}\frac{17}{6}
\end{align}
Let $\theta = \tan^{-1}\frac{17}{6}$. From \eqref{eq:CBSE12-2019-sin_arctan}:
\begin{align}
	\sin\brak{\tan^{-1}\frac{17}{6}} = \frac{17}{\sqrt{17^2 + 6^2}} = \frac{17}{\sqrt{289 + 36}} = \frac{17}{\sqrt{325}} = \frac{17}{5\sqrt{13}} = \frac{17\sqrt{13}}{65}
\end{align}


\item Using properties of determinants, show that
\begin{align}
	\mydet{
		3a & -a+b & -a+c \\
		-b+a & 3b & -b+c \\
		-c+a & -c+b & 3c
	} = 3\brak{a+b+c}\brak{ab+bc+ca}
\end{align}
\hfill (12~2019)

\solution
Let $\Delta$ denote the given determinant. Applying the column operation $C_1 \to C_1 + C_2 + C_3$ in accordance with \eqref{eq:CBSE12-2019-det_shear_invariance}:
\begin{align}
	\Delta = \mydet{
		a+b+c & -a+b & -a+c \\
		a+b+c & 3b & -b+c \\
		a+b+c & -c+b & 3c
	}
\end{align}
Factoring out $\brak{a+b+c}$ from column $C_1$:
\begin{align}
	\Delta = \brak{a+b+c} \mydet{
		1 & -a+b & -a+c \\
		1 & 3b & -b+c \\
		1 & -c+b & 3c
	}
\end{align}
Applying row operations $R_2 \to R_2 - R_1$ and $R_3 \to R_3 - R_1$:
\begin{align}
	\Delta = \brak{a+b+c} \mydet{
		1 & -a+b & -a+c \\
		0 & a+2b & a-b \\
		0 & a-c & a+2c
	}
\end{align}
Expanding along column $C_1$:
\begin{align}
	\Delta &= \brak{a+b+c} \mydet{a+2b & a-b \\ a-c & a+2c} \\
	&= \brak{a+b+c} \sbrak{\brak{a+2b}\brak{a+2c} - \brak{a-b}\brak{a-c}} \\
	&= \brak{a+b+c} \sbrak{\brak{a^2 + 2ac + 2ab + 4bc} - \brak{a^2 - ac - ab + bc}} \\
	&= \brak{a+b+c}\brak{3ab + 3bc + 3ca} \\
	&= 3\brak{a+b+c}\brak{ab + bc + ca}
\end{align}


\item Check whether the relation $R$ defined on the set $A = \{1, 2, 3, 4, 5, 6\}$ by $R = \{(a, b) : b = a + 1\}$ is reflexive, symmetric or transitive.
\hfill (12~2019)

\solution

Listing all ordered pairs satisfying $b = a + 1$ for $a, b \in A$:
\begin{align}
	R = \{(1, 2), \, (2, 3), \, (3, 4), \, (4, 5), \, (5, 6)\}
\end{align}

\begin{itemize}
	\item \textbf{Reflexivity:} A relation is reflexive if $(a, a) \in R$ for all $a \in A$. Since $(1, 1) \notin R$, the relation is \textbf{not reflexive}.
	\item \textbf{Symmetry:} A relation is symmetric if $(a, b) \in R \implies (b, a) \in R$. Since $(1, 2) \in R$ but $(2, 1) \notin R$, the relation is \textbf{not symmetric}.
	\item \textbf{Transitivity:} A relation is transitive if $(a, b) \in R$ and $(b, c) \in R \implies (a, c) \in R$. Since $(1, 2) \in R$ and $(2, 3) \in R$ but $(1, 3) \notin R$, the relation is \textbf{not transitive}.
\end{itemize}
Thus, $R$ is neither reflexive, nor symmetric, nor transitive.

\item Let $f:\mathbb{N}\to Y$ be a function defined as $f(x)=4x+3$, where $Y=\{y\in\mathbb{N}:y=4x+3\text{ for some }x\in\mathbb{N}\}$. Show that $f$ is invertible. Find its inverse.
\hfill (12~2019)

\solution
From \eqref{eq:CBSE12-2019-injective}, to prove $f$ is injective (one-one), let $f(x_1) = f(x_2)$ for $x_1, x_2 \in \mathbb{N}$:
\begin{align}
	4x_1 + 3 = 4x_2 + 3 \implies 4x_1 = 4x_2 \implies x_1 = x_2
\end{align}
Thus, $f$ is injective.
From \eqref{eq:CBSE12-2019-surjective}, for any $y \in Y$, by definition $y = 4x + 3$ for some $x \in \mathbb{N}$. Solving for $x$:
\begin{align}
	x = \frac{y - 3}{4} \in \mathbb{N}
\end{align}
Since every $y \in Y$ has a pre-image in $\mathbb{N}$, $f$ is surjective (onto).
Since $f$ is both injective and surjective, it is bijective and hence invertible.
Defining the inverse function $f^{-1}: Y \to \mathbb{N}$:
\begin{align}
	f^{-1}(y) = \frac{y - 3}{4}
\end{align}

\item If $x\sqrt{1+y}+y\sqrt{1+x}=0$, where $x \neq y$, prove that $\dfrac{dy}{dx} = -\dfrac{1}{\brak{x+1}^2}$.
\hfill (12~2019)

\solution
Rearranging the given equation:
\begin{align}
	x\sqrt{1+y} &= -y\sqrt{1+x}
\end{align}
Squaring both sides:
\begin{align}
	x^2\brak{1+y} &= y^2\brak{1+x} \\
	x^2 + x^2y &= y^2 + xy^2 \\
	x^2 - y^2 + x^2y - xy^2 &= 0 \\
	\brak{x - y}\brak{x + y} + xy\brak{x - y} &= 0 \\
	\brak{x - y}\sbrak{x + y + xy} &= 0
\end{align}
Since $x \neq y$, we have $x - y \neq 0$, which implies:
\begin{align}
	x + y + xy = 0 \implies y\brak{1 + x} = -x \implies y = -\frac{x}{1 + x}
\end{align}
Rewriting $y$ to facilitate differentiation:
\begin{align}
	y = -\frac{\brak{x + 1} - 1}{x + 1} = -\brak{1 - \frac{1}{x+1}} = -1 + \brak{x+1}^{-1}
\end{align}
Differentiating with respect to $x$:
\begin{align}
	\frac{dy}{dx} = \frac{d}{dx}\sbrak{-1 + \brak{x+1}^{-1}} = -\brak{x+1}^{-2}\cdot 1 = -\frac{1}{\brak{x+1}^2}
\end{align}


\item If $\brak{\cos x}^y = \brak{\sin y}^x$, find $\dfrac{dy}{dx}$.
\hfill (12~2019)

\solution
Taking the natural logarithm on both sides using \eqref{eq:CBSE12-2019-log_diff}:
\begin{align}
	\ln\left(\brak{\cos x}^y\right) &= \ln\left(\brak{\sin y}^x\right) \\
	y \ln\brak{\cos x} &= x \ln\brak{\sin y}
\end{align}
Differentiating both sides implicitly with respect to $x$ using \eqref{eq:CBSE12-2019-product_log_diff}:
\begin{align}
	\frac{dy}{dx}\ln\brak{\cos x} + y\cdot\frac{-\sin x}{\cos x} &= 1\cdot\ln\brak{\sin y} + x\cdot\frac{\cos y}{\sin y}\frac{dy}{dx} \\
	\frac{dy}{dx}\ln\brak{\cos x} - y\tan x &= \ln\brak{\sin y} + x\cot y\frac{dy}{dx}
\end{align}
Grouping the derivative terms:
\begin{align}
	\frac{dy}{dx}\sbrak{\ln\brak{\cos x} - x\cot y} &= \ln\brak{\sin y} + y\tan x
\end{align}
Solving for $\dfrac{dy}{dx}$:
\begin{align}
	\frac{dy}{dx} = \frac{\ln\brak{\sin y} + y\tan x}{\ln\brak{\cos x} - x\cot y}
\end{align}



\item Prove that $\displaystyle\int_0^a f(x)\,dx = \int_0^a f(a-x)\,dx$, and hence evaluate $\displaystyle\int_0^{\pi/2}\frac{x}{\sin x+\cos x}\,dx$.
\hfill (12~2019)

\solution
To prove King's property \eqref{eq:CBSE12-2019-kings_property}, let $t = a - x \implies dx = -dt$. When $x = 0$, $t = a$; when $x = a$, $t = 0$:
\begin{align}
	\int_0^a f(x)\,dx = \int_a^0 f(a - t)\brak{-dt} = \int_0^a f(a - t)\,dt = \int_0^a f(a - x)\,dx
\end{align}
Now, evaluating the given definite integral:
\begin{align}
	I = \int_0^{\pi/2} \frac{x}{\sin x + \cos x}\,dx \label{eq:EE26BTECH11223-12-q23-integral}
\end{align}
Applying property \eqref{eq:CBSE12-2019-kings_property}:
\begin{align}
	I &= \int_0^{\pi/2} \frac{\frac{\pi}{2} - x}{\sin\left(\frac{\pi}{2} - x\right) + \cos\left(\frac{\pi}{2} - x\right)}\,dx \\
	&= \int_0^{\pi/2} \frac{\frac{\pi}{2} - x}{\cos x + \sin x}\,dx \label{eq:EE26BTECH11223-10-q23-king}
\end{align}
Adding \eqref{eq:EE26BTECH11223-12-q23-integral} and \eqref{eq:EE26BTECH11223-10-q23-king}:
\begin{align}
	2I &= \int_0^{\pi/2} \frac{x + \frac{\pi}{2} - x}{\sin x + \cos x}\,dx = \frac{\pi}{2} \int_0^{\pi/2} \frac{1}{\sin x + \cos x}\,dx \\
	I &= \frac{\pi}{4} \int_0^{\pi/2} \frac{1}{\sqrt{2}\cos\left(x - \frac{\pi}{4}\right)}\,dx = \frac{\pi}{4\sqrt{2}} \int_0^{\pi/2} \sec\left(x - \frac{\pi}{4}\right)\,dx
\end{align}
Substituting $u = x - \frac{\pi}{4}$:
\begin{align}
	I &= \frac{\pi}{4\sqrt{2}} \int_{-\pi/4}^{\pi/4} \sec u\,du = \frac{\pi}{4\sqrt{2}} \cdot 2\ln\left|\sec\frac{\pi}{4} + \tan\frac{\pi}{4}\right| \\
	&= \frac{\pi}{2\sqrt{2}}\ln\left(\sqrt{2} + 1\right)
\end{align}


\item Find $\displaystyle\int \frac{x^2+x+1}{\brak{x+2}\brak{x^2+1}}\,dx$.
\hfill (12~2019)

\solution
Using partial fraction decomposition from \eqref{eq:CBSE12-2019-partial_fractions}:
\begin{align}
	\frac{x^2 + x + 1}{\brak{x + 2}\brak{x^2 + 1}} = \frac{A}{x + 2} + \frac{Bx + C}{x^2 + 1}
\end{align}
Multiplying by $\brak{x + 2}\brak{x^2 + 1}$:
\begin{align}
	x^2 + x + 1 &= A\brak{x^2 + 1} + \brak{Bx + C}\brak{x + 2} \\
	&= \brak{A + B}x^2 + \brak{2B + C}x + \brak{A + 2C}
\end{align}
Comparing coefficients of like powers of $x$:
\begin{align}
	A + B &= 1 \\
	2B + C &= 1 \\
	A + 2C &= 1
\end{align}
Solving this linear system yields $A = \frac{3}{5}$, $B = \frac{2}{5}$, and $C = \frac{1}{5}$. Thus:
\begin{align}
	\frac{x^2 + x + 1}{\brak{x + 2}\brak{x^2 + 1}} = \frac{3}{5}\brak{\frac{1}{x + 2}} + \frac{1}{5}\brak{\frac{2x}{x^2 + 1}} + \frac{1}{5}\brak{\frac{1}{x^2 + 1}}
\end{align}
Integrating term-by-term using \eqref{eq:CBSE12-2019-int_linear_denom}--\eqref{eq:CBSE12-2019-int_arctan}:
\begin{align}
	I &= \frac{3}{5}\int \frac{1}{x + 2}\,dx + \frac{1}{5}\int \frac{2x}{x^2 + 1}\,dx + \frac{1}{5}\int \frac{1}{x^2 + 1}\,dx \\
	&= \frac{3}{5}\ln\abs{x + 2} + \frac{1}{5}\ln\brak{x^2 + 1} + \frac{1}{5}\tan^{-1}x + C
\end{align}


\item Solve the differential equation $x\dfrac{dy}{dx} = y - x\tan\left(\frac{y}{x}\right)$.
\hfill (12~2019)

\solution
Rewriting in the homogeneous form \eqref{eq:CBSE12-2019-homogeneous_ode}:
\begin{align}
	\frac{dy}{dx} = \frac{y}{x} - \tan\left(\frac{y}{x}\right)
\end{align}
Substituting $y = vx$ and $\dfrac{dy}{dx} = v + x\dfrac{dv}{dx}$ in accordance with \eqref{eq:CBSE12-2019-homog_separable}:
\begin{align}
	v + x\frac{dv}{dx} = v - \tan v \implies x\frac{dv}{dx} = -\tan v
\end{align}
Separating variables:
\begin{align}
	\frac{dv}{\tan v} = -\frac{dx}{x} \implies \cot v\,dv = -\frac{dx}{x}
\end{align}
Integrating both sides:
\begin{align}
	\ln\abs{\sin v} = -\ln\abs{x} + \ln C_1 \implies \ln\abs{x\sin v} = \ln C_1 \implies x\sin v = C
\end{align}
Substituting back $v = \frac{y}{x}$:
\begin{align}
	x\sin\left(\frac{y}{x}\right) = C
\end{align}



\item Solve the differential equation $\dfrac{dy}{dx} = -\dfrac{x + y\cos x}{1 + \sin x}$.
\hfill (12~2019)

\solution
Rearranging into the standard linear differential form \eqref{eq:CBSE12-2019-linear_ode}:
\begin{align}
	\frac{dy}{dx} + \left(\frac{\cos x}{1 + \sin x}\right)y = -\frac{x}{1 + \sin x}
\end{align}
Here $P(x) = \frac{\cos x}{1 + \sin x}$ and $Q(x) = -\frac{x}{1 + \sin x}$.
Finding the integrating factor via \eqref{eq:CBSE12-2019-linear_ode_sol}:
\begin{align}
	\text{I.F.} = e^{\int \frac{\cos x}{1 + \sin x}\,dx} = e^{\ln\brak{1 + \sin x}} = 1 + \sin x
\end{align}
Multiplying through and integrating:
\begin{align}
	y\brak{1 + \sin x} &= \int \left(-\frac{x}{1 + \sin x}\right)\brak{1 + \sin x}\,dx + C \\
	&= \int (-x)\,dx + C = -\frac{x^2}{2} + C
\end{align}
Multiplying by $2$:
\begin{align}
	2y\brak{1 + \sin x} + x^2 = C'
\end{align}



\item The scalar product of the vector $\vec{a} = \vec{i} + \vec{j} + \vec{k}$ with a unit vector along the sum of the vectors $\vec{b} = 2\vec{i} + 4\vec{j} - 5\vec{k}$ and $\vec{c} = \lambda\vec{i} + 2\vec{j} + 3\vec{k}$ is equal to $1$. Find the value of $\lambda$ and hence find the unit vector along $\vec{b} + \vec{c}$.
\hfill (12~2019)

\solution
Summing vectors $\vec{b}$ and $\vec{c}$:
\begin{align}
	\vec{b} + \vec{c} = \myvec{2 + \lambda \\ 6 \\ -2}
\end{align}
The unit vector $\vec{u}$ along $\vec{b} + \vec{c}$ is given by \eqref{eq:CBSE12-2019-unit_vector}:
\begin{align}
	\vec{u} = \frac{\vec{b} + \vec{c}}{\norm{\vec{b} + \vec{c}}} = \frac{\myvec{\lambda + 2 \\ 6 \\ -2}}{\sqrt{\brak{\lambda + 2}^2 + 36 + 4}} = \frac{\myvec{\lambda + 2 \\ 6 \\ -2}}{\sqrt{\brak{\lambda + 2}^2 + 40}}
\end{align}
Given that $\vec{a} \cdot \vec{u} = 1$:
\begin{align}
	\myvec{1 & 1 & 1} \cdot \frac{\myvec{\lambda + 2 \\ 6 \\ -2}}{\sqrt{\brak{\lambda + 2}^2 + 40}} = 1 \implies \frac{\brak{\lambda + 2} + 6 - 2}{\sqrt{\brak{\lambda + 2}^2 + 40}} = 1 \implies \frac{\lambda + 6}{\sqrt{\brak{\lambda + 2}^2 + 40}} = 1
\end{align}
Squaring both sides ($\lambda + 6 > 0$):
\begin{align}
	\brak{\lambda + 6}^2 = \brak{\lambda + 2}^2 + 40 \implies \lambda^2 + 12\lambda + 36 = \lambda^2 + 4\lambda + 4 + 40 \implies 8\lambda = 8 \implies \lambda = 1
\end{align}
Substituting $\lambda = 1$:
\begin{align}
	\vec{b} + \vec{c} = \myvec{3 \\ 6 \\ -2}, \quad \norm{\vec{b} + \vec{c}} = \sqrt{9 + 36 + 4} = \sqrt{49} = 7
\end{align}
Thus, the unit vector along $\vec{b} + \vec{c}$ is:
\begin{align}
	\vec{u} = \frac{1}{7}\myvec{3 \\ 6 \\ -2} = \frac{3}{7}\vec{i} + \frac{6}{7}\vec{j} - \frac{2}{7}\vec{k}
\end{align}


\item If $\brak{a + bx}e^{y/x} = x$, then prove that $x^3\dfrac{d^2y}{dx^2} = \left(x\dfrac{dy}{dx} - y\right)^2$.
\hfill (12~2019)

\solution
Taking the natural logarithm on both sides:
\begin{align}
	\ln\brak{a + bx} + \frac{y}{x} = \ln x \implies \frac{y}{x} = \ln x - \ln\brak{a + bx}
\end{align}
Differentiating with respect to $x$:
\begin{align}
	\frac{x\frac{dy}{dx} - y}{x^2} = \frac{1}{x} - \frac{b}{a + bx} = \frac{a + bx - bx}{x\brak{a + bx}} = \frac{a}{x\brak{a + bx}}
\end{align}
Multiplying by $x^2$:
\begin{align}
	x\frac{dy}{dx} - y = \frac{ax}{a + bx} \label{eq:EE26BTECH11223-12-q28-rel1}
\end{align}
Differentiating \eqref{eq:EE26BTECH11223-12-q28-rel1} with respect to $x$:
\begin{align}
	\frac{dy}{dx} + x\frac{d^2y}{dx^2} - \frac{dy}{dx} &= \frac{\brak{a + bx}\brak{a} - \brak{ax}\brak{b}}{\brak{a + bx}^2} \\
	x\frac{d^2y}{dx^2} &= \frac{a^2}{\brak{a + bx}^2}
\end{align}
Multiplying both sides by $x^2$:
\begin{align}
	x^3\frac{d^2y}{dx^2} = \frac{a^2 x^2}{\brak{a + bx}^2} = \left(\frac{ax}{a + bx}\right)^2
\end{align}
Substituting \eqref{eq:EE26BTECH11223-12-q28-rel1}:
\begin{align}
	x^3\frac{d^2y}{dx^2} = \left(x\frac{dy}{dx} - y\right)^2
\end{align}
The following code verifies the differential identity:
\begin{lstlisting}[frame=single, basicstyle=\ttfamily\small]
	codes/12/q28.py
\end{lstlisting}

\item The volume of a cube is increasing at the rate of $8\text{ cm}^3/\text{s}$. How fast is the surface area increasing when the length of its edge is $12\text{ cm}$?
\hfill (12~2019)

\solution
Let $x$ be the edge length of the cube. From \eqref{eq:CBSE12-2019-cube_rates}:
\begin{align}
	V = x^3 \implies \frac{dV}{dt} = 3x^2 \frac{dx}{dt}
\end{align}
Given $\dfrac{dV}{dt} = 8\text{ cm}^3/\text{s}$:
\begin{align}
	8 = 3x^2 \frac{dx}{dt} \implies \frac{dx}{dt} = \frac{8}{3x^2}
\end{align}
Surface area $S = 6x^2$. From \eqref{eq:CBSE12-2019-cube_rates}, the rate of change of surface area is:
\begin{align}
	\frac{dS}{dt} = 12x \frac{dx}{dt} = 12x \left(\frac{8}{3x^2}\right) = \frac{32}{x}
\end{align}
Evaluating at $x = 12\text{ cm}$:
\begin{align}
	\left.\frac{dS}{dt}\right|_{x=12} = \frac{32}{12} = \frac{8}{3}\text{ cm}^2/\text{s}
\end{align}


\item Find the Cartesian and vector equations of the plane passing through the points $A(2,5,-3)$, $B(-2,-3,5)$, and $C(5,3,-3)$.
\hfill (12~2019)

\solution
Let the position vectors of the points be
\begin{align}
	\vec{a} = \myvec{2 \\ 5 \\ -3}, \quad \vec{b} = \myvec{-2 \\ -3 \\ 5}, \quad \vec{c} = \myvec{5 \\ 3 \\ -3}
\end{align}
Evaluating displacement vectors lying in the plane:
\begin{align}
	\vec{AB} &= \vec{b} - \vec{a} = \myvec{-4 \\ -8 \\ 8} \\
	\vec{AC} &= \vec{c} - \vec{a} = \myvec{3 \\ -2 \\ 0}
\end{align}
From \eqref{eq:CBSE12-2019-plane_vector_normal}, the normal vector $\vec{n}$ is obtained by cross product:
\begin{align}
	\vec{n} = \vec{AB} \times \vec{AC} = \mydet{\vec{i} & \vec{j} & \vec{k} \\ -4 & -8 & 8 \\ 3 & -2 & 0} = \vec{i}\brak{0 - \brak{-16}} - \vec{j}\brak{0 - 24} + \vec{k}\brak{8 - \brak{-24}} = 16\vec{i} + 24\vec{j} + 32\vec{k}
\end{align}
Scaling by $\frac{1}{8}$ yields the simplified normal vector:
\begin{align}
	\vec{n}_0 = \myvec{2 \\ 3 \\ 4}
\end{align}
From \eqref{eq:CBSE12-2019-plane_vector_normal}, evaluating $\vec{a} \cdot \vec{n}_0$:
\begin{align}
	\myvec{2 & 5 & -3}\myvec{2 \\ 3 \\ 4} = 4 + 15 - 12 = 7
\end{align}
Thus, the vector equation of the plane is:
\begin{align}
	\vec{r} \cdot \myvec{2 \\ 3 \\ 4} = 7
\end{align}
Substituting $\vec{r} = \myvec{x \\ y \\ z}$, the Cartesian equation is:
\begin{align}
	2x + 3y + 4z = 7
\end{align}
The following code verifies the plane equations:
\begin{lstlisting}[frame=single, basicstyle=\ttfamily\small]
	codes/12/q30.py
\end{lstlisting}

\item Find the point on the curve $y^2 = 4x$ which is nearest to the point $(2, -8)$.
\hfill (12~2019)

\solution
Let $P\brak{x, y}$ be a point on the parabola $y^2 = 4x$, so its coordinates can be expressed parametrically as $\myvec{y^2/4 \\ y}$.
Let the fixed point be $\vec{A} = \myvec{2 \\ -8}$.
The square of the Euclidean distance $Z$ between $P$ and $A$ from \eqref{eq:CBSE12-2019-extrema_first} is:
\begin{align}
	Z = \brak{\frac{y^2}{4} - 2}^2 + \brak{y - \brak{-8}}^2 = \brak{\frac{y^2}{4} - 2}^2 + \brak{y + 8}^2
\end{align}
Differentiating $Z$ with respect to $y$:
\begin{align}
	\frac{dZ}{dy} &= 2\brak{\frac{y^2}{4} - 2}\brak{\frac{2y}{4}} + 2\brak{y + 8} \\
	&= \frac{y^3}{4} - 2y + 2y + 16 = \frac{y^3}{4} + 16
\end{align}
Setting $\frac{dZ}{dy} = 0$ for stationary points:
\begin{align}
	\frac{y^3}{4} + 16 = 0 \implies y^3 = -64 \implies y = -4
\end{align}
Applying the second derivative test \eqref{eq:CBSE12-2019-extrema_second}:
\begin{align}
	\frac{d^2Z}{dy^2} = \frac{3y^2}{4}
\end{align}
Evaluating at $y = -4$:
\begin{align}
	\left.\frac{d^2Z}{dy^2}\right|_{y=-4} = \frac{3\brak{-4}^2}{4} = \frac{3\brak{16}}{4} = 12 > 0
\end{align}
Since $\frac{d^2Z}{dy^2} > 0$, $Z$ attains its minimum at $y = -4$.
Substituting $y = -4$ into $x = \frac{y^2}{4}$:
\begin{align}
	x = \frac{\brak{-4}^2}{4} = \frac{16}{4} = 4
\end{align}
Therefore, the point on the curve nearest to $(2, -8)$ is $(4, -4)$.
The following code verifies the optimization:
\begin{lstlisting}[frame=single, basicstyle=\ttfamily\small]
	codes/12/q31.py
\end{lstlisting}

\item Find $\displaystyle\int_1^3 \brak{x^2 + 2 + e^{2x}}\,dx$ as the limit of sums.
\hfill (12~2019)

\solution
From the Riemann sum definition \eqref{eq:CBSE12-2019-riemann_limit}:
\begin{align}
	\int_a^b f(x)\,dx = \lim_{h \to 0} h \sum_{r=1}^n f\brak{a + rh}
\end{align}
where $a = 1, b = 3 \implies nh = b - a = 2$, and $f(x) = x^2 + 2 + e^{2x}$.
Evaluating $f\brak{1 + rh}$:
\begin{align}
	f\brak{1 + rh} = \brak{1 + rh}^2 + 2 + e^{2\brak{1 + rh}} = 3 + 2rh + r^2h^2 + e^2\brak{e^{2h}}^r
\end{align}
Using standard summation formulas \eqref{eq:CBSE12-2019-sum_r}, \eqref{eq:CBSE12-2019-sum_r2}, and \eqref{eq:CBSE12-2019-sum_geom}:
\begin{align}
	h S_n &= 3\brak{nh} + 2h^2\left[\frac{n\brak{n+1}}{2}\right] + h^3\left[\frac{n\brak{n+1}\brak{2n+1}}{6}\right] + e^2 e^{2h} \left(\frac{e^{2nh} - 1}{e^{2h} - 1}\right) h \\
	&= 3\brak{nh} + \brak{nh}\brak{nh + h} + \frac{\brak{nh}\brak{nh + h}\brak{2nh + h}}{6} + e^2\brak{e^{2nh} - 1}e^{2h} \left(\frac{2h}{e^{2h} - 1}\right)\frac{1}{2}
\end{align}
Taking the limit as $h \to 0$ using \eqref{eq:CBSE12-2019-exp_limit}:
\begin{align}
	\int_1^3 \brak{x^2 + 2 + e^{2x}}\,dx &= 3\brak{2} + 2\brak{2} + \frac{2\brak{2}\brak{4}}{6} + \brak{e^6 - e^2}\brak{1}\left(\frac{1}{2}\cdot 1\right) \\
	&= 6 + 4 + \frac{8}{3} + \frac{e^6 - e^2}{2} = \frac{38}{3} + \frac{e^6 - e^2}{2}
\end{align}


\item Using integration, find the area of the triangular region whose sides have the equations $y = 2x + 1$, $y = 3x + 1$, and $x = 4$.
\hfill (12~2019)

\solution
Solving the lines simultaneously:
\begin{itemize}
	\item Intersection of $y = 2x + 1$ and $y = 3x + 1 \implies x = 0, y = 1$ (Vertex $A(0, 1)$).
	\item Intersection of $y = 2x + 1$ and $x = 4 \implies x = 4, y = 9$ (Vertex $B(4, 9)$).
	\item Intersection of $y = 3x + 1$ and $x = 4 \implies x = 4, y = 13$ (Vertex $C(4, 13)$).
\end{itemize}
From \eqref{eq:CBSE12-2019-area_curves}, integrating between $x = 0$ and $x = 4$:
\begin{align}
	\text{Area} &= \int_0^4 \sbrak{\brak{3x + 1} - \brak{2x + 1}}\,dx = \int_0^4 x\,dx \\
	&= \left.\frac{x^2}{2}\right|_0^4 = \frac{4^2}{2} = 8\text{ square units}
\end{align}
The following code verifies the area:
\begin{lstlisting}[frame=single, basicstyle=\ttfamily\small]
	codes/12/q33.py
\end{lstlisting}

\item If $\vec{A} = \begin{bmatrix} 1 & 3 & 4 \\ 2 & 1 & 2 \\ 5 & 1 & 1 \end{bmatrix}$, find $\vec{A}^{-1}$. Hence solve the system of equations:
\begin{align}
	x + 3y + 4z &= 8 \\
	2x + y + 2z &= 5 \\
	5x + y + z &= 7
\end{align}
\hfill (12~2019)

\solution
Evaluating the determinant of $\vec{A}$:
\begin{align}
	\mydet{\vec{A}} = 1\brak{1 - 2} - 3\brak{2 - 10} + 4\brak{2 - 5} = 1\brak{-1} - 3\brak{-8} + 4\brak{-3} = -1 + 24 - 12 = 11
\end{align}
Using \eqref{eq:CBSE12-2019-inv_adjugate}, computing the cofactor matrix and its transpose (adjugate):
\begin{align}
	\operatorname{adj}(\vec{A}) = \begin{bmatrix}
		-1 & 1 & 2 \\
		8 & -19 & 6 \\
		-3 & 14 & -5
	\end{bmatrix}
\end{align}
Thus, $\vec{A}^{-1}$ is given by:
\begin{align}
	\vec{A}^{-1} = \frac{1}{11}\begin{bmatrix}
		-1 & 1 & 2 \\
		8 & -19 & 6 \\
		-3 & 14 & -5
	\end{bmatrix}
\end{align}
Writing the system of equations in matrix form $\vec{A}\vec{X} = \vec{B}$ where $\vec{X} = \myvec{x \\ y \\ z}$ and $\vec{B} = \myvec{8 \\ 5 \\ 7}$:
\begin{align}
	\vec{X} = \vec{A}^{-1}\vec{B} = \frac{1}{11}\begin{bmatrix}
		-1 & 1 & 2 \\
		8 & -19 & 6 \\
		-3 & 14 & -5
	\end{bmatrix}\myvec{8 \\ 5 \\ 7} = \frac{1}{11}\myvec{-8 + 5 + 14 \\ 64 - 95 + 42 \\ -24 + 70 - 35} = \frac{1}{11}\myvec{11 \\ 11 \\ 11} = \myvec{1 \\ 1 \\ 1}
\end{align}
Therefore, $x = 1$, $y = 1$, and $z = 1$.


\item Find the inverse of the matrix $\vec{A} = \begin{bmatrix} 2 & 0 & -1 \\ 5 & 1 & 0 \\ 0 & 1 & 3 \end{bmatrix}$ using elementary transformations.
\hfill (12~2019)

\solution
From \eqref{eq:CBSE12-2019-gauss_jordan}, setting up the augmented matrix $\sbrak{\vec{A} \mid \vec{I}}$ and applying Gauss-Jordan row elimination:
\begin{align}
	\sbrak{\vec{A} \mid \vec{I}} = \left[\begin{array}{ccc|ccc} 2 & 0 & -1 & 1 & 0 & 0 \\ 5 & 1 & 0 & 0 & 1 & 0 \\ 0 & 1 & 3 & 0 & 0 & 1 \end{array}\right]
\end{align}
Performing elementary row transformations $R_2 \to R_2 - 2R_1$, row swaps, and elimination steps reduces the left block to the identity matrix $\vec{I}$, yielding:
\begin{align}
	\vec{A}^{-1} = \begin{bmatrix}
		3 & -1 & 1 \\
		-15 & 6 & -5 \\
		5 & -2 & 2
	\end{bmatrix}
\end{align}


\item A company produces two types of goods, A and B, that require gold and silver. Each unit of type A requires $3\text{ g}$ of silver and $1\text{ g}$ of gold while that of type B requires $1\text{ g}$ of silver and $2\text{ g}$ of gold. The company can use at most $9\text{ g}$ of silver and $8\text{ g}$ of gold. If each unit of type A brings a profit of Rs.~40 and that of type B Rs.~50, find the number of units of each type that the company should produce to maximize profit. Formulate the above LPP and solve it graphically and find the maximum profit.
\hfill (12~2019)

\solution
Let $x$ and $y$ be the units of type A and type B goods produced.
From \eqref{eq:CBSE12-2019-lpp_obj} and \eqref{eq:CBSE12-2019-lpp_constraints}, the LPP formulation is:
\begin{align}
	\text{Maximize } & Z = 40x + 50y \\
	\text{subject to } & 3x + y \le 9 \\
	& x + 2y \le 8 \\
	& x \ge 0, \; y \ge 0
\end{align}
The feasible region is the polygon $OABC$ bounded by the extreme corner points:
\begin{align}
	O(0, 0), \quad A(3, 0), \quad B(2, 3), \quad C(0, 4)
\end{align}
Evaluating the objective function $Z$ at each corner point:
\begin{itemize}
	\item At $O(0, 0): Z = 0$
	\item At $A(3, 0): Z = 40\brak{3} + 50\brak{0} = 120$
	\item At $B(2, 3): Z = 40\brak{2} + 50\brak{3} = 80 + 150 = 230$
	\item At $C(0, 4): Z = 40\brak{0} + 50\brak{4} = 200$
\end{itemize}
By the Corner-Point Theorem, the maximum profit is Rs.~$230$, achieved by producing $2$ units of type A and $3$ units of type B.
The following code verifies the graphical optimization:
\begin{lstlisting}[frame=single, basicstyle=\ttfamily\small]
	codes/12/q36.py
\end{lstlisting}

\item There are three coins. One is a two-headed coin, another is a biased coin that comes up heads $75\%$ of the time and the third is an unbiased coin. One of the three coins is chosen at random and tossed. If it shows heads, what is the probability that it is the two-headed coin?
\hfill (12~2019)

\solution
Let $E_1, E_2, E_3$ be the events of choosing the two-headed coin, the biased coin, and the unbiased coin respectively.
\begin{align}
	\pr{E_1} = \pr{E_2} = \pr{E_3} = \frac{1}{3}
\end{align}
Let $A$ be the event that the tossed coin shows heads. The conditional probabilities are:
\begin{align}
	\pr{A \mid E_1} = 1, \quad \pr{A \mid E_2} = \frac{75}{100} = \frac{3}{4}, \quad \pr{A \mid E_3} = \frac{1}{2}
\end{align}
Using Bayes' theorem \eqref{eq:CBSE12-2019-bayes_theorem}:
\begin{align}
	\pr{E_1 \mid A} &= \frac{\pr{E_1}\pr{A \mid E_1}}{\pr{E_1}\pr{A \mid E_1} + \pr{E_2}\pr{A \mid E_2} + \pr{E_3}\pr{A \mid E_3}} \\
	&= \frac{\frac{1}{3}\brak{1}}{\frac{1}{3}\brak{1} + \frac{1}{3}\brak{\frac{3}{4}} + \frac{1}{3}\brak{\frac{1}{2}}} = \frac{1}{1 + \frac{3}{4} + \frac{1}{2}} = \frac{1}{\frac{9}{4}} = \frac{4}{9}
\end{align}


\item Find the equation of the line passing through $(2,-1,2)$ and $(5,3,4)$ and of the plane passing through $(2,0,3)$, $(1,1,5)$ and $(3,2,4)$. Also, find their point of intersection.
\hfill (12~2019)

\solution
The direction ratios of the line passing through $(2, -1, 2)$ and $(5, 3, 4)$ are $3, 4, 2$. Its parametric Cartesian equation from \eqref{eq:CBSE12-2019-line_cart} is:
\begin{align}
	\frac{x - 2}{3} = \frac{y + 1}{4} = \frac{z - 2}{2} = \lambda \implies x = 3\lambda + 2, \; y = 4\lambda - 1, \; z = 2\lambda + 2
\end{align}
The plane passing through $(2, 0, 3)$, $(1, 1, 5)$, and $(3, 2, 4)$ has normal $\vec{n} = \myvec{-1 \\ 1 \\ 2} \times \myvec{1 \\ 2 \\ 1} = \myvec{-3 \\ 3 \\ -3}$, yielding normal vector $\vec{n}_0 = \myvec{1 \\ -1 \\ 1}$.
From \eqref{eq:CBSE12-2019-plane_three_points}, the Cartesian equation of the plane is:
\begin{align}
	x - y + z = 5
\end{align}
From \eqref{eq:CBSE12-2019-line_plane_intersection}, substituting the line parameterization into the plane equation:
\begin{align}
	\brak{3\lambda + 2} - \brak{4\lambda - 1} + \brak{2\lambda + 2} = 5 \implies \lambda + 5 = 5 \implies \lambda = 0
\end{align}
Substituting $\lambda = 0$ into the line equations yields the point of intersection:
\begin{align}
	\vec{P} = \myvec{2 \\ -1 \\ 2}
\end{align}
The following code verifies the intersection:
\begin{lstlisting}[frame=single, basicstyle=\ttfamily\small]
	codes/12/q38.py
\end{lstlisting}

\end{enumerate}


